% ECONOMICS_PROBLEM: {"id": "E52-572", "title": "Quantitative tightening transmission", "jel_code": "E52", "source_id": "OP-0246"}
% !TeX program = xelatex
\documentclass[11pt,a4paper]{article}
\usepackage[margin=23mm]{geometry}
\usepackage{fontspec}
\setmainfont{Latin Modern Roman}
\usepackage{amsmath,amssymb,mathtools}
\usepackage{xcolor,enumitem,booktabs,longtable,array,xurl,hyperref}
\definecolor{muted}{HTML}{53636C}
\hypersetup{colorlinks=true,urlcolor=blue,linkcolor=blue}
\setlength{\parindent}{0pt}
\setlength{\parskip}{5pt}
\newcommand{\R}{\mathbb R}
\newcommand{\E}{\mathbb E}
\newcommand{\N}{\mathbb N}
\newcommand{\ind}{\mathbf 1}
\newcommand{\Var}{\operatorname{Var}}
\newcommand{\Cov}{\operatorname{Cov}}
\newcommand{\supp}{\operatorname{supp}}
\DeclareMathOperator*{\argmin}{arg\,min}
\DeclareMathOperator*{\argmax}{arg\,max}
\newcommand{\entrymeta}[3]{{\sffamily\small\color{muted}\textbf{\texttt{#1}}\quad|\quad #2\par #3\par}}
\newcommand{\Problemref}[1]{\texttt{#1}}
\newcommand{\JELref}[2]{\texttt{#2} (JEL \texttt{#1})}
\begin{document}
\section*{Quantitative tightening transmission}
\textbf{Problem ID:} \texttt{E52-572}\quad \textbf{JEL:} \texttt{E52}\par
% BEGIN ECONOMICS STATEMENT
\label{problem:OP-0246}
{\sffamily\small\color{muted}Original subject group: Inflation and monetary policy\par}
\label{definitions:problem:30007527}
\textbf{Audit classification:} Published research agenda. Normalisation is an agenda with existing firm-response evidence; conditional QE/QT symmetry is editorial. Verification concerns the cited version, not first priority or proof that this exact editorial specification remains unsolved on October 8, 2026.
\subsection*{Definitions and precise research target}
\paragraph{Editorial mathematical specification.} Let \(b_t\) be an unanticipated central-bank purchase of duration-adjusted government securities, with purchases positive and sales or runoff surprises negative. Fix a maturity segment and define predetermined state \(s_t\) to include reserve abundance, intermediary balance-sheet capacity, and market stress. Let \(F_{j,t+h}(b)\) be firm \(j\)'s potential financing cost, debt maturity, or investment at horizon \(h\), with other policy surprises held fixed. For \(a>0\), define
\[
A_h(a,s)=E[F_{j,t+h}(a)-F_{j,t+h}(0)\mid s_t=s]
+E[F_{j,t+h}(-a)-F_{j,t+h}(0)\mid s_t=s].
\]
Quantitative easing and tightening are symmetric at this scale and state if \(A_h(a,s)=0\). Identify the function \(A_h\), with confidence regions, from announcement surprises that separate duration supply, anticipated runoff, and policy-rate information. Match or reweight the support of easing and tightening episodes; if states do not overlap, report bounds or model-based extrapolation rather than an unconditional symmetry conclusion. Distinguish financing responses from real activity and report heterogeneity by predetermined credit quality. The target is scale- and state-specific symmetry, including possible liquidity thresholds, not a claim that every purchase and sale should mechanically have equal and opposite effects.

Let duration-adjusted purchases be the sum of changes in security holdings times a fixed duration weight; define the unit and maturity bin before analysis. Runoff is a reduction in holdings, but anticipated runoff is not itself an innovation. Define the potential outcomes through a joint intervention on balance-sheet purchases and the policy-rate surprise, with the latter fixed at zero. The asymmetry function is the sum of two differences in the same financing-outcome units and state, so zero tests oddness around the zero-intervention baseline; it does not test equality across historical QE and QT regimes. Identification requires conditional shock exogeneity and support at \(-a,0,a\) or an explicitly stated structural extrapolation. Conditioning on reserve abundance must use its preannouncement value.
\subsection*{Variants and assumption changes}
The following are \emph{editorial variants}, not additional published open conjectures.
\begin{enumerate}
\item \emph{Liquidity threshold.} Define a predetermined reserve-to-bank-assets ratio \(R\) and estimate the symmetry function separately above and below a fixed scarcity threshold \(R_0\). The threshold itself is a separate estimand if not fixed in advance.
\item \emph{Maturity and credit quality.} Keep total duration supplied constant while varying the purchased maturity segment. Compare high- and low-credit-quality firms with common policy-shock exposure; financing and investment effects remain separate outcomes.
\end{enumerate}
\subsection*{References and source audit}
\begin{itemize}
\item European Central Bank. \emph{Monetary policy, strategy and implementation}. Undated. Locator: Policy normalisation, selected areas of focus. Primary text checked. \url{https://www.ecb.europa.eu/pub/economic-research/research_agenda/monetary_policy/html/index.en.html}.
\item Egemen Eren; Denis Gorea; Daojing Zhai. \emph{How do quantitative easing and tightening affect firms?}. 2025. Locator: BIS WP1286 primary metadata, Focus/\allowbreak Contribution/\allowbreak Findings/\allowbreak Abstract; September 2, 2025. Primary text checked. \url{https://www.bis.org/publications/working-paper-1286-how-do-quantitative-easing-and-tightening-affect-firms}.
\end{itemize}
Normalisation is an agenda with existing firm-response evidence; conditional QE/QT symmetry is editorial.
\begin{enumerate}\raggedright
\item\hypertarget{definitions:source:30007527:1}{} European Central Bank. Undated / date unverified. \emph{Monetary policy, strategy and implementation}. Policy normalisation, selected areas of focus.
\url{https://www.ecb.europa.eu/pub/economic-research/research_agenda/monetary_policy/html/index.en.html}.
\item\hypertarget{definitions:source:30007527:2}{} Egemen Eren; Denis Gorea; Daojing Zhai. 2025. \emph{How do quantitative easing and tightening affect firms?}. BIS WP1286 primary metadata, Focus/\allowbreak Contribution/\allowbreak Findings/\allowbreak Abstract; September2,2025.
\url{https://www.bis.org/publications/working-paper-1286-how-do-quantitative-easing-and-tightening-affect-firms}.
\end{enumerate}

\subsection*{Common conventions: Source mathematical conventions}

These conventions apply to each entry unless explicitly replaced.

\textbf{Models and identification.} A primitive model \(m\) specifies all probability laws, feasible choices, information, equilibrium and measurement rules used by its target. The admissible class \(\mathcal M\) is fixed independently of outcomes. If \(P_O(m)\) is its observable probability law and \(\tau(m)\) the target, its identified set at observable law \(P\) is
\[
 \mathcal I_\tau(P)=\{\tau(m):m\in\mathcal M,\ P_O(m)=P\}.
\]
Point identification means that this set is a singleton. An empty set signals incompatibility of data and assumptions. Coordinate bounds need not describe the joint set. Bounds are sharp only if every retained target is attainable or arbitrarily approachable by compatible models. Finite bounds require explicit restrictions on unobserved quantities; unrestricted confounding, hidden wealth, extrapolation or arbitrary equilibrium selection can yield uninformative sets.

\textbf{Causal interventions.} A potential outcome \(Y_i(p)\) is the outcome under a complete policy regime \(p\), including timing, financing, information and market spillovers. The target population and units are fixed before treatment. With interference, use the complete assignment vector or a justified exposure mapping; individual no-interference notation alone cannot identify an economy-wide intervention. A difference of mean potential outcomes does not require the cross-world joint distribution. Conditioning on post-treatment migration, survival, enrollment or employment changes the estimand and needs separate justification.

\textbf{Statistical statements.} An estimator, confidence set, forecast or test specifies a sampling law, model class and loss/coverage criterion. Uniform \((1-\alpha)\)-coverage means \(\inf_{m\in\mathcal M}\Pr_m\{\tau(m)\in C_n\}\geq1-\alpha\), or an explicitly asymptotic version. The whole parameter space trivially has coverage, so informativeness requires a stated width, risk or power objective. A forecasting improvement is relative to a named benchmark, information set and evaluation population.

\textbf{Equilibrium and optimization.} An equilibrium comprises feasible plans, optimality, beliefs where needed, budget constraints and all market/resource clearing. The primitives or a stated benchmark define each correspondence \(\mathcal E(p;m)\). Nonemptiness is assumed or proved, never inferred just from compact policy sets. A selected-equilibrium target uses a declared selection rule; a robust target quantifies over all permitted equilibria. Use suprema or infima unless attainment is established. An \(\varepsilon\)-optimal policy is feasible and within \(\varepsilon\) of the finite optimum value. Report an empty feasible set separately.

\textbf{Welfare and accounting.} Normative utility, interpersonal normalization, social weights, numeraire, discounting and the use of public receipts are primitives. Behavioral utility need not coincide with the welfare criterion. Household welfare includes ownership income and rents once; adding profits again to compensating variation generally double-counts them. A monetary equivalent is defined by an explicit transfer and price convention, with monotonicity and range conditions. Average effects, mechanism contributions, transfers and welfare losses are different targets. Infinite sums require convergence and dynamic feasible plans require terminal or no-Ponzi conditions.

\textbf{Source versions.} A working-paper date, revision date and journal publication year are distinguished. A DOI for a journal version is not automatically the DOI of an earlier manuscript. Locators refer to the stated version and printed pages where specified. A metadata-only check does not verify a claim about a full-text conclusion. Every entry's source subsection narrows the provenance claim accordingly.
% END ECONOMICS STATEMENT
\end{document}
